\chapter{Conclusões}
\label{chap:conclusao}

Este trabalho apresentou a implementação de uma plataforma web para organizar informações de produtores, pontos de coleta, amostras, análises e arquivos no contexto do NAPI Abelhas. A solução materializa parte central da proposta do TCC 1: associa a amostra à sua origem, mantém análises vinculadas e disponibiliza o envio de documentos por armazenamento de objetos. A interface em Next.js e a API em NestJS, com PostgreSQL, Prisma, Clerk e MinIO, constituem os artefatos técnicos entregues.

A comparação com o planejamento mostrou mudanças de tecnologia e de acesso. O frontend passou de Vite para Next.js, a API foi estruturada em NestJS sobre Express, e o perfil único previsto deu lugar a organizações com papéis de administrador e membro. Também evidenciou requisitos ainda não demonstrados: filtros avançados no servidor, registro do local físico da amostra, imagem obrigatória no cadastro e uso efetivo de Redis como cache. Assim, o resultado deve ser compreendido como uma implementação com limites identificados, e não como comprovação de atendimento integral ao escopo inicialmente proposto.

Como continuidade, recomenda-se completar os requisitos pendentes, ampliar os testes automatizados dos fluxos de negócio e executar testes de ponta a ponta em ambiente configurado com dados fictícios. Uma avaliação posterior com pesquisadores e técnicos poderá verificar a adequação dos formulários, da navegação e das consultas às rotinas reais do NAPI Abelhas. Somente com essa avaliação e com medições específicas será possível afirmar efeitos da plataforma sobre produtividade, desempenho ou qualidade das informações.
